% Part:sets-functions-relations
% Chapter: size-of-sets
% Section: zig-zag

\documentclass[../../../include/open-logic-section]{subfiles}

\begin{document}

\olfileid{sfr}{siz}{zigzag}

\olsection{کانتور دا زگ زیگ طریقہ}

\begin{explain}
اسی کجھ ”سوکھیاں“ گنتی وار فہرستاں پہلے ای ویکھ چکے آں۔ ہُن اسی
کجھ ذرا اوکھا ویکھاں گے۔ قدرتی عدداں دیاں جوڑیاں دے سیٹ اُتے غور
کرو\oliflabeldef{sfr:set:pai:sec}{، جیہنوں اسی \olref[set][pai]{sec}
وچ ایویں تعریف کیتا سی:}{، جیہڑا ایویں تعریف اے:}
\[
\Nat \times \Nat = \Setabs{\tuple{n,m}}{n,m \in \Nat}
\]
اسی ایہناں ترتیب وار جوڑیاں نوں اک \emph{جدول} وچ ایویں سجا سکدے آں:
\[
\begin{array}{ c | c | c | c | c | c}
& \mathbf 0 & \mathbf 1 & \mathbf 2 & \mathbf 3 & \dots \\
\hline
\mathbf 0 & \tuple{0,0} & \tuple{0,1} & \tuple{0,2} & \tuple{0,3} & \dots \\
\hline
\mathbf 1 & \tuple{1,0} & \tuple{1,1} & \tuple{1,2} & \tuple{1,3} & \dots \\
\hline
\mathbf 2 & \tuple{2,0} & \tuple{2,1} & \tuple{2,2} & \tuple{2,3} & \dots \\
\hline
\mathbf 3 & \tuple{3,0} & \tuple{3,1} & \tuple{3,2} & \tuple{3,3} & \dots \\
\hline
\vdots & \vdots & \vdots & \vdots & \vdots & \ddots\\
\end{array}
\]
صاف اے کہ $\Nat \times \Nat$ دی ہر ترتیب وار جوڑی جدول وچ ٹھیک اک
واری آوے گی۔ خاص طور تے، $\tuple{n,m}$، $n$ویں قطار تے $m$ویں کالم
وچ آوے گی۔ پر اسی ایہو جہے جدول دے عنصراں نوں ”اک جہتی“ فہرست وچ
کِویں سجاواں؟ ہیٹھلے جدول دا نمونہ ایہ کرن دا اک طریقہ وکھاؤندا اے
(بھاویں ہور کئی طریقے وی موجود نیں):
\[
\begin{array}{ c | c | c | c | c | c | c}
& \mathbf 0 & \mathbf 1 & \mathbf 2 & \mathbf 3 & \mathbf 4 &\dots \\
\hline
\mathbf 0 & 0  & 1& 3 & 6& 10 &\ldots \\
\hline
\mathbf 1 &2 & 4& 7 & 11 & \dots &\ldots \\
\hline
\mathbf 2 & 5 & 8 & 12 & \ldots & \dots&\ldots \\
\hline
\mathbf 3 & 9 & 13 & \ldots & \ldots & \dots & \ldots \\
\hline
\mathbf 4 & 14 & \ldots & \ldots & \ldots & \dots & \ldots \\
\hline
\vdots & \vdots & \vdots & \vdots & \vdots&\ldots & \ddots\\
\end{array}
\]\noindent
ایس نمونے نوں \emph{کانتور دا زگ زیگ طریقہ} آکھیا جاندا اے۔ ایہ
$\Nat \times \Nat$ دی گنتی وار فہرست ایویں بناؤندا اے:
\[
\tuple{0,0}, \tuple{0,1}, \tuple{1,0}, \tuple{0,2}, \tuple{1,1},
\tuple{2,0}, \tuple{0,3}, \tuple{1,2}, \tuple{2,1}, \tuple{3,0}, \dots
\]
تے ایس توں اگلی گل قائم ہوندی اے:
\end{explain}

\begin{prop}\ollabel{natsquaredenumerable}
$\Nat \times \Nat$، !!{enumerable} اے۔
\end{prop}

\begin{proof}
فنکشن $f \colon \Nat \to \Nat\times\Nat$، ہر $k \in \Nat$ نوں
اوہ تِک $\tuple{n,m} \in \Nat \times \Nat$ دیندا اے جیہدے لئی $k$،
کانتور دے زگ زیگ جدول دی $n$ویں قطار تے $m$ویں کالم دی قدر اے۔
\end{proof}

\begin{explain}
ایس تکنیک نوں بڑے سوہنے ڈھنگ نال عام وی کیتا جا سکدا اے۔ مثال لئی،
اسی قدرتی عدداں دے ترتیب وار تِکاں دے سیٹ دی گنتی وار فہرست بناؤن
لئی اینوں ورت سکدے آں، یعنی:
\[
\Nat \times \Nat \times \Nat = \Setabs{\tuple{n,m,k}}{n,m,k \in \Nat}
\]
اسی $\Nat \times \Nat \times \Nat$ نوں $\Nat \times \Nat$ تے
$\Nat$ دا کارتیسی حاصل ضرب سمجھدے آں، یعنی
\[
\Nat^3 = (\Nat \times \Nat) \times \Nat =
\Setabs{\tuple{\tuple{n,m},k}}{n, m, k
  \in \Nat }
\]
تے ایس طرح اسی $\Nat^3$ دی گنتی وار فہرست اک جدول نال بنا سکدے آں:
اک محور اُتے $\Nat$ دی گنتی وار فہرست دے لیبل لا دیو، تے دوجے محور
اُتے $\Nat^2$ دی گنتی وار فہرست دے:
\[
\begin{array}{ c | c | c | c | c | c}
& \mathbf 0 & \mathbf 1 & \mathbf 2 & \mathbf 3 & \dots \\
\hline
\mathbf{\tuple{0,0}} & \tuple{0,0,0} & \tuple{0,0,1} & \tuple{0,0,2} & \tuple{0,0,3} & \dots \\
\hline
\mathbf{\tuple{0,1}} & \tuple{0,1,0} & \tuple{0,1,1} & \tuple{0,1,2} & \tuple{0,1,3} & \dots \\
\hline
\mathbf{\tuple{1,0}} & \tuple{1,0,0} & \tuple{1,0,1} & \tuple{1,0,2} & \tuple{1,0,3} & \dots \\
\hline
\mathbf{\tuple{0,2}} & \tuple{0,2,0} & \tuple{0,2,1} & \tuple{0,2,2} & \tuple{0,2,3} & \dots\\
\hline
\vdots & \vdots & \vdots & \vdots & \vdots & \ddots \\
\end{array}
\]
سو کانتور دے زگ زیگ طریقے ورگا طریقہ ورت کے اسی اوسے طرح $\Nat^3$
دی گنتی وار فہرست حاصل کر سکدے آں۔ تے اسی ایویں اگے ودھدے ہوئے،
$\Nat^n$ دیاں گنتی وار فہرستاں حاصل کر سکدے آں، ہر قدرتی عدد $n$ لئی۔
سو ساڈے کول اگلا نتیجہ اے:
\end{explain}

\begin{prop}
$\Nat^n$، ہر $n \in \Nat$ لئی، !!{enumerable} اے۔
\end{prop}

\begin{prob}\label{sfr:siz:zigzag:prob:posint-n}
وکھاؤ کہ $(\PosInt)^n$، ہر $n \in \Nat$ لئی، !!{enumerable} اے۔
\end{prob}

\begin{prob}\label{sfr:siz:zigzag:prob:posint-star}
  وکھاؤ کہ $(\PosInt)^*$، !!{enumerable} اے۔ تسیں
  \cref{sfr:siz:zigzag:prob:posint-n} منّ سکدے او۔
\end{prob}

\end{document}
